\section{Introduction}\label{sec:introduction}

As university research infrastructure grows, the volume and diversity of specialised equipment catalogued in institutional repositories increases proportionally, making efficient discovery a non-trivial challenge. At the University of Groningen, the database currently indexes over 8,000 instruments; however, its keyword-based search interface requires researchers to know the precise name or category of the equipment they seek. When a researcher can only describe an experimental need, rather than a specific instrument, the existing system fails to surface relevant results, directly impeding experimental efficiency.

To address this limitation, this thesis proposes an advanced search system designed to act as the foundational retrieval component for a Retrieval-Augmented Generation (RAG) pipeline. RAG augments a language model's generative capacity with a retrieval component that dynamically queries an external knowledge base at inference time \citep{gao2023retrieval, lewis2020retrieval}. This approach is particularly suited to the domain because the equipment catalogue is updated continuously, and retrieval-based systems can incorporate new documents without retraining, a critical advantage over fine-tuning-based alternatives. Furthermore, RAG confines responses to grounded, retrievable evidence, reducing hallucination risk and ensuring factual accuracy. Rather than requiring researchers to memorise equipment names, the proposed system allows them to describe their experimental context in natural language; the retrieval component then identifies and ranks the most appropriate instruments, together with metadata such as access permissions and physical location.

The scope of this thesis is deliberately bounded: it focuses exclusively on the retrieval component of the RAG pipeline, systematically comparing and evaluating multiple retrieval strategies: sparse (BM25), dense (bi-encoder), and hybrid, rather than examining the downstream generation stage. The central research contribution is a rigorous analysis that aims to answer the following research question: \textbf{How do different retrieval strategies and embedding techniques impact the accuracy of equipment recommendations in RAG-based search system?}

